\section{Problem Definition and Stakeholder Context}
\label{sec:problem}

\subsection{Two problems, one decision chain}

The exploratory analysis of Section~\ref{sec:eda} left one finding standing
above the rest: a late delivery raises the probability of a 1--2 star review
from 9.2\% to 54.0\%, a 5.9-fold increase monotone in how late the parcel is.
The project turns that sentence into a pair of predictions placed at two points
in an order's life (Table~\ref{tab:problems}).

\begin{table}[htbp]
\centering
\caption{The two problems. They share one feature pipeline, one model-family
budget and one dataset; what separates them is \emph{when} the model runs.}
\label{tab:problems}
\footnotesize
\begin{tabular}{@{}p{1.9cm}p{6.2cm}p{6.2cm}@{}}
\toprule
 & \textbf{Problem 1 --- delivery performance} & \textbf{Problem 2 --- satisfaction risk} \\
\midrule
Question & How long will this order take, and will it miss its promise? &
Will this delivered order draw a 1--2 star review? \\[2pt]
Moment & \code{at\_checkout} --- the order has just been placed &
\code{at\_delivery} --- the parcel has arrived, the review is not yet written \\[2pt]
Stakeholder & Operations and logistics & Customer experience \\[2pt]
Task & \textbf{1A regression}: lead time in days, on MAE.
\textbf{1B classification}: delivered after the promised date, on PR-AUC against
its no-skill floor &
\textbf{Classification}: \code{review\_score} $\le 2$, on PR-AUC against its
floor and balanced accuracy \\[2pt]
Action & Re-quote the promise, warn the buyer, expedite freight &
Reach out before the review lands \\
\bottomrule
\end{tabular}
\end{table}

The brief permits two problems or one framed twice; we take two, because they
are halves of one decision chain: operations acts before the parcel ships,
customer experience after it arrives and before the customer writes anything,
and Problem 1's output is, in substance, Problem 2's most important input.
Problem 1 is the brief's \emph{dual framing}: one underlying quantity, two
outputs, so the late-order flag \emph{is} the regression's answer compared
against the promised date; five classifiers trained directly on the binary
label run alongside as a control arm (Section~\ref{sec:results-late}).

\subsection{How the stakeholders act on the output}
\label{sec:levers}

A prediction is only worth producing if someone changes a decision because of
it --- the lesson of the prescriptive delivery framework reviewed in
Section~\ref{sec:lit}, where the savings came from routing machine-learned
delivery outcomes into a scheduling decision \citep{kandula2021prescriptive}.
Table~\ref{tab:levers} states what each team does with each output and what
decides whether the action pays, and it shapes the evaluation: the cheap
operations levers apply to every order and are judged on the regression's
accuracy; the expensive one, and all of customer experience's work, are rationed
by capacity and are judged on what the top-$K$ of a ranked list delivers
(Section~\ref{sec:results-capacity}).

\begin{table}[htbp]
\centering
\caption{What each team does with the predictions, and what decides whether
the action pays. Costs are the stated assumptions of Section~\ref{sec:cost}.}
\label{tab:levers}
\footnotesize
\begin{tabular}{@{}p{2.0cm}p{3.6cm}p{1.35cm}p{6.9cm}@{}}
\toprule
\textbf{Who, when} & \textbf{Lever} & \textbf{Cost} & \textbf{Trigger, and the evidence behind it} \\
\midrule
Operations, at checkout & Re-quote the promised date with the model's estimate & none & every order; MAE 3.90 days against the promise's 9.79 (Section~\ref{sec:results-reg}) \\[2pt]
 & Pre-warn the buyer & near zero & $p_{\text{late}} \ge 0.11$ catches 82.6\% of the orders that will miss \\[2pt]
 & Expedite freight & R\$12 & only where precision clears 0.30, which it never does (Section~\ref{sec:results-capacity}); manual review of the top of the list \\[2pt]
 & Monitor sellers, routes, carrier & analyst time & track-record features and the weekly late series (Figure~\ref{fig:lateness}) as early warning \\[2pt]
\midrule
Customer experience, at delivery & Contact the customer before the review is written & R\$5 & the day's deliveries ranked by $p_{\text{low}}$, worked top-down to capacity $K$; at 5\% of deliveries, two contacts in three reach a would-be 1--2 star reviewer \\
\bottomrule
\end{tabular}
\end{table}

\subsection{Targets, and what the review score measures}

All three targets come from the provided tables alone. \emph{Lead time} is the
time from purchase (\code{order\_purchase\_\allowbreak timestamp}) to arrival
(\code{order\_delivered\_\allowbreak customer\_date}), in days, defined for the
97.0\% of orders that reached the customer. \emph{Late} is delivery strictly after
\code{order\_estimated\_delivery\_date}, which carries no time component, so
lateness is accurate to about a day. \emph{Low review} is \code{review\_score}
$\le 2$, defined for the 99.2\% carrying a review. Undefined rows are \code{NaN}
and excluded, never imputed to a negative label: the event has not happened
rather than gone missing, and a silent customer is not a satisfied one.

The target of Problem 2 is a platform rating, not the \emph{perceived quality}
of the survey literature, and the distinction --- which the Phase 1 evaluation
asked for --- matters for what the model claims. Perceived service quality is
measured by instruments such as E-S-QUAL \citep{parasuraman2005esqual} and
e-LSQ \citep{rao2011elsq}: self-reported and retrospective. A star rating is a
behavioural trace left by the customers who choose to write one, and its
J-shaped distribution --- 57.8\% five stars here, with a second mode at one
star --- is the signature of acquisition and under-reporting biases
\citep{hu2009jshaped, hu2017selfselection}; a rating also moves with the
product, its price and the seller's communication
\citep{dellarocas2003digitization}. Problem 2 therefore predicts the
\emph{rating}, the observable the customer-experience team is measured on, and
claims nothing about latent satisfaction; the consequence, that part of the
residual is unobserved product and seller effects, is carried into
Section~\ref{sec:limitations}.

\subsection{Success criteria, registered in advance}
\label{sec:bands}

Expected ranges for all six headline metrics were written into
\code{config.EXPECTED\_BANDS} \emph{before} any model was tuned, from the
Phase 1 results and untuned scoping probes (Appendix~\ref{app:bands}). They are
the targets the work is judged against and they are a leak alarm: here the
signature of leakage is not a subtle improvement but an implausible one
(Section~\ref{sec:leakage}), so a result landing far \emph{above} its band is
investigated before it is reported. All six landed inside; none above. The
ceiling was stated as plainly as the target: the largest swings in lateness are
shocks to carrier capacity and to the platform's own promising policy
(Section~\ref{sec:eda-lateness}), none of which is an attribute of an order, so
a late-delivery classifier reporting 95\% accuracy here would be the
majority-class rule or a leak.
